% Source-aligned English for printed p.180.
\SourcePage{180}

\LayerTarget{tohoku:corollary:3.9.1:3}
\noindent\textsc{Corollary 3.} \textit{Suppose that the preliminary condition
of the preceding corollary holds, and suppose in addition that the coverings
$\mathbf R\in\mathfrak R(A)$ have nerves of dimension at most $n$. Then the
equivalent conditions of the preceding corollary in fact imply
$H^i(A,F)=0$ for $A\in\mathcal S$ and every $i>0$.}

\LayerTarget{tohoku:proof:corollary-3.9.1-3}
Indeed, one automatically has $H^i(\mathbf R,F)=0$ for
$\mathbf R\in\mathfrak R(A)$ and $i>n$, and hence for every $i$. That is,
condition $A(m)$ holds for every $m$, so it is enough to apply the corollary
with $m$ large.

\LayerTarget{tohoku:proof:theorem-3.9.1:induction-on-n}
\noindent\textsc{Proof of Theorem 3.9.1.} We proceed by induction on the
integer $n$, the assertion being clear when $n=0$. Suppose then that
$n\geq1$ and that the theorem has been proved for integers $n'<n$. The
induction hypothesis first gives $H^i(A,F)=\{0\}$ for
$1\leq i\leq n-1$. It therefore remains to prove that
$c^n\in H^n(A,F)$ is zero, under the condition that there is a finite subset
$L$ of $\mathfrak R(A)$ such that the restriction of $c^n$ to every
$B\in\mathbf R^L$ is zero. Let $k$ be the number of elements of $L$; we
argue by induction on $k$. The conclusion is trivial if $k=0$; let us prove
it for $k=1$.

\LayerTarget{tohoku:proof:theorem-3.9.1:base-k-one}
By hypothesis there is a $\mathbf R\in\mathfrak R(A)$ such that the
restriction of $c^n$ to every $B\in\mathbf R$ is zero. Since the members
$B\in\mathbf R$ belong to $\mathcal S$, we have $H^i(B,F)=\{0\}$ for
$1\leq i\leq n-1$, so the term $E_2^{p,q}$ of the Leray spectral sequence
relative to $\mathbf R$ is zero for $1\leq q\leq n-1$. In the usual way this
gives an exact sequence

\LayerTarget{tohoku:eq:theorem-3.9.1-edge-sequence}
\[
H^n(\mathbf R,F)\longrightarrow H^n(A,F)
\longrightarrow H^0(\mathbf R,H^n(F))\longrightarrow\cdots.
\]

The first term is zero by $A(n)$, so the second homomorphism is injective. By
hypothesis the image of $c^n$ under this homomorphism is zero; hence $c^n$ is
zero.

\LayerTarget{tohoku:proof:theorem-3.9.1:induction-on-k}
Now suppose that $k\geq2$ and that the conclusion has been proved for
$k'<k$. Let $L=(\mathbf R^1,\ldots,\mathbf R^k)$. In view of what has just
been proved, it is enough to show that the restriction $c_B^n$ of $c^n$ to
every $B\in\mathbf R^1$ is zero. For $i=2,\ldots,k$, the trace
$\mathbf R_B^i$ of $\mathbf R^i$ on $B$ belongs to $\mathfrak R(B)$.
Moreover, by the hypothesis on $c^n$, the restriction of $c_B^n$ to every
member of the intersection of the coverings $\mathbf R_B^i$, for
$2\leq i\leq k$, is zero. Applying the induction hypothesis for $k'=k-1$ to
$c_B^n$ and $B$, we obtain $c_B^n=0$, completing the proof of the theorem.

\LayerTarget{tohoku:proposition:3.9.2}
\noindent\textsc{Proposition 3.9.2.} \textit{The preliminary hypothesis of
Corollary 2 holds in the following case: $X$ is quasi-compact, the members
$A\in\mathcal S$ are closed, $X\in\mathcal S$, the coverings
$\mathbf R\in\mathfrak R(X)$ are finite and, for every $A\in\mathcal S$ and
every $\mathbf R\in\mathfrak R(X)$, the trace $\mathbf R_A$ of $\mathbf R$ on
$A$ belongs to $\mathfrak R(A)$; moreover, for any two distinct points $x,y$
of $X$ there is a $\mathbf R\in\mathfrak R(X)$ no member of which contains
both $x$ and $y$.}

\LayerTarget{tohoku:definition:quasi-compact}
(We shall say that a space is quasi-compact if it satisfies the open-cover
axiom for compact spaces, without necessarily being Hausdorff.)

\LayerTarget{tohoku:construction:cover-stars-and-neighborhoods}
\noindent\textsc{Proof.} For a covering $\mathbf R$ of $X$ and a point
$x\in X$, let $E_x(\mathbf R)$ (the star of $\mathbf R$ at $x$) be the union
of the members $A\in\mathbf R$ that contain $x$, and let $O_x(\mathbf R)$ be
the complement of the union of the members $A\in\mathbf R$ that do not
contain $x$.

\LayerTarget{tohoku:eq:cover-star-inclusions}
Thus $x\in O_x(\mathbf R)\subset E_x(\mathbf R)$; moreover, for every
$y\in O_x(\mathbf R)$ one has $E_y(\mathbf R)\subset E_x(\mathbf R)$.

\LayerTarget{tohoku:claim:closed-cover-star-neighborhood}
If $\mathbf R$ is a finite covering of $X$ by closed sets, then
$O_x(\mathbf R)$ is an \textit{open neighborhood of $x$}.

\LayerTarget{tohoku:proof:proposition-3.9.2:finite-intersection:p1}
Under the conditions of the proposition, let $\mathcal U$ be an open covering
of $X$, let $x\in X$, and choose $U_x\in\mathcal U$ with $x\in U_x$. By
hypothesis, the intersection of the $E_x(\mathbf R)$ as $\mathbf R$ ranges
over $\mathfrak R(X)$ is reduced to $x$. Quasi-compactness therefore gives a
finite subset $L_x$ of $\mathfrak R(X)$ such that the intersection of the
$E_x(\mathbf R)$ for $\mathbf R\in L_x$, i.e.
$E_x(\mathbf R^{L_x})$, is contained in $U_x$.

\TohokuContinuation{181}
